\begin{definition}
  Let $G$ be any group. 
  A $G$-torsor is an inhabited set $X$ with a right action 
  $(\cdot): X \times G \to X$ such that the induced map 
  $\langle 1 , (\cdot)\rangle : X \times G \to X \times X$ 
  is a bijection. 

  We denote $BG$ for the type of $G$-torsors.
\end{definition}
The intuition for $G$-torsors is that they are groups who have forgotten their unit element. 
\begin{example}
  $G$ with it's multiplication is itself a $G$-torsor. 
\end{example}
\begin{definition}
  Let $X$ be a $G$-torsor and $x,y:X$. We denote $x^{-1}\cdot y:G$ as the unique element $g$ satisying $x\cdot g = y$. 
\end{definition}
\begin{remark}
  Note that $(x^{-1} \cdot y) \cdot (y^{-1} \cdot z) = x^{-1} \cdot z$. 
\end{remark}
\begin{remark}
  For any $G$-torsor $X$, the map $X \to (X \to G)$ mapping $x,y$ to $x^{-1}\cdot y$ induces an equivalence $X \simeq (X \simeq G)$.
\end{remark}
\begin{corollary}
  The loop space of the pointed type $(BG,G)$ is equivalent to $G$. This justifies that we call $BG$ the \textbf{delooping} of $G$. 
\end{corollary}
\begin{remark}
Under the identification $X \simeq (X \simeq G)$ applied to $X = G$, the reflexive path is identified with the identity element $0$ of $G$. In this paper, we freely switch between the identification described above. For example, we identify $\alpha = 0$ with $\alpha$.
\end{remark}


